% =====================================================================================
% REVISIÓN 1: solo los «Temas tentativos del marco teórico».
% REVISIÓN 2: el Marco teórico completo (cuando lo desarrolles, sustituye o convierte los temas tentativos en subsecciones).
% Escribe tu texto FUERA de los recuadros \begin{guia} ... \end{guia}.
% =====================================================================================

\section{Marco teórico} \label{sec:marco_teorico} % audit:req=marco_teorico

\begin{guia}[Guía: Marco teórico (Revisión 2)]
\textbf{¿Para qué sirve?} Reúne los conceptos, teorías, tecnologías, métodos y normas que se necesitan para entender y
justificar tu propuesta. No es un glosario ni un resumen de Internet: cada tema se explica \emph{en función de tu proyecto}.

\textbf{Debe responder (por cada tema):}
\begin{itemize}[nosep,leftmargin=*]
  \item ¿Qué es y cómo funciona? (definición con fuente citada y, si ayuda, un esquema o ecuación).
  \item ¿Para qué se usa en mi proyecto y por qué lo elegí frente a las alternativas?
  \item ¿Qué ventajas, limitaciones o buenas prácticas debo conocer?
\end{itemize}
Además: ¿qué metodología de desarrollo o investigación usarás y por qué? ¿qué normas o estándares aplican (calidad,
seguridad, protección de datos)?

\textbf{Organización:} una \verb|\subsection| por tema, de lo general a lo particular.

\textbf{Extensión mínima:} 400 palabras de desarrollo (sin contar los temas tentativos de la Revisión 1).

\textbf{Fuentes:} libros, artículos y documentación oficial. Cada afirmación importante lleva \verb|\cite{clave}|.

\textbf{Errores frecuentes:} temas sin relación con el proyecto; copiar y pegar; fuentes que son solo URLs de blogs; no
citar; temas que luego no se usan en el desarrollo.

\textbf{Cómo se revisa:} pertinencia, profundidad suficiente, calidad de las fuentes y que cada tema se vincule con la
propuesta.
\end{guia}

% >>> REVISIÓN 2: ESCRIBE AQUÍ el marco teórico (introducción breve y una \subsection por tema) <<<


\subsection{Temas tentativos del marco teórico} \label{sec:temas_tentativos} % audit:req=temas_marco_teorico

\begin{guia}[Guía: Temas tentativos del marco teórico (Revisión 1)]
\textbf{¿Para qué sirve?} En la primera revisión no se pide el marco teórico completo, sino la lista de los temas que
vas a desarrollar. Así el director puede confirmar que vas bien encaminado.

\textbf{Debe responder:} ¿cuáles son los temas (de 5 a 10) que necesito para sustentar el proyecto? Para cada uno, en una
o dos oraciones: ¿qué es y por qué lo necesito?

\textbf{Formato:} una lista (\verb|\begin{enumerate}| o \verb|itemize|), un tema por elemento.

\textbf{Extensión mínima:} 25 palabras. En la Revisión 2 conviertes cada tema en una subsección desarrollada.

\textbf{Cómo se revisa:} que los temas cubran lo necesario para el proyecto y no incluyan relleno.
\end{guia}

% >>> REVISIÓN 1: ESCRIBE AQUÍ la lista de temas tentativos <<<

\begin{enumerate}

\item \textbf{Grandes Modelos de Lenguaje (LLM).} Los Grandes Modelos de Lenguaje son sistemas de inteligencia artificial entrenados con grandes cantidades de texto para procesar y generar lenguaje natural. Este tema es necesario para comprender la base de los modelos que serán analizados en el proyecto y el contexto en el que surgen las propuestas de LLM de 1 bit.

\item \textbf{Arquitectura Transformer.} La arquitectura Transformer es una estructura utilizada en numerosos modelos de lenguaje modernos y permite procesar relaciones entre elementos de una secuencia mediante mecanismos de atención. Este tema es necesario para comprender la base arquitectónica sobre la que se desarrollan distintos LLM y cómo se relaciona con los modelos de baja precisión estudiados.

\item \textbf{Cuantización de modelos de lenguaje.} La cuantización es una técnica que reduce la precisión numérica utilizada para representar los parámetros de un modelo. Este tema es necesario para comprender cómo se disminuyen los requerimientos de memoria y procesamiento y cómo se relaciona este proceso con los modelos de lenguaje de 1 bit.

\item \textbf{Modelos de lenguaje de 1 bit.} Los modelos de lenguaje de 1 bit utilizan representaciones de muy baja precisión en sus parámetros con el objetivo de reducir los recursos computacionales necesarios para su almacenamiento y ejecución. Este tema es central para el proyecto porque constituye el objeto principal del estudio comparativo.

\item \textbf{Métricas de evaluación de modelos de lenguaje.} Las métricas de evaluación permiten medir distintos aspectos del comportamiento y desempeño de un modelo de lenguaje. Este tema es necesario para identificar criterios que permitan comparar de manera objetiva los modelos seleccionados durante el desarrollo del proyecto.

\item \textbf{Eficiencia computacional y uso de recursos en LLM.} La eficiencia computacional analiza los recursos necesarios para entrenar o ejecutar un modelo, como memoria, almacenamiento y capacidad de procesamiento. Este tema es necesario para valorar la viabilidad técnica de los modelos de 1 bit y comparar sus requerimientos bajo condiciones similares.

\end{enumerate}
